% Part: lambda-calculus
% Chapter: syntax
% Section: terms

\documentclass[../../../include/open-logic-section]{subfiles}

\begin{document}

\olfileid[pt]{lam}{syn}{trm}
\olsection{Termos}

Os termos do cálculo lambda são construídos indutivamente a partir de um
suprimento infinito de variáveis $\Obj{v_0}$, $\Obj{v_1}$, \dots, do símbolo
~``$\lambd$'' e de parênteses. Usaremos $x$, $y$, $z$, \dots{} para
designar variáveis, e $M$, $N$, $P$, \dots{} para designar termos.

\begin{defn}[Termos] \ollabel{defn:term}
O conjunto de \emph{termos} do cálculo lambda é definido indutivamente por:
\begin{enumerate}
  \item \ollabel{defn:term-var} Se $x$ é uma variável, então $x$ é um
    termo.
  \item \ollabel{defn:term-abs} Se $x$ é uma variável e $M$ é um
    termo, então $(\lambd[x][M])$ é um termo.
  \item \ollabel{defn:term-app} Se $M$ e $N$ são termos, então
    $(MN)$ é um termo.
\end{enumerate}
\end{defn}

Se um termo $(\lambd[x][M])$ é formado de acordo com
\olref{defn:term-abs}, dizemos que ele é resultado de uma
\emph{abstração}, e o $x$ em $\lambd[x]$ é chamado de
\emph{!!{parameter}}. Um termo $(MN)$ formado de acordo com
\olref{defn:term-app} é resultado de uma \emph{aplicação}.

Os termos definidos acima têm todos os parênteses explicitados. Isso pode se
tornar bastante incômodo, como demonstra o termo
$(\lambd[x][((\lambd[x][x])(\lambd[x][(xx)]))])$. Introduziremos convenções
para evitar parênteses. A definição oficial, porém, facilita determinar como
um termo é construído segundo \olref{defn:term}. Por exemplo, o último passo
na formação do termo
$(\lambd[x][((\lambd[x][x])(\lambd[x][(xx)]))])$ deve ser uma abstração
cujo !!{parameter} é~$x$. Ele resulta da abstração aplicada ao termo
$((\lambd[x][x])(\lambd[x][(xx)]))$, que é uma aplicação de dois termos.
Cada um desses dois termos é resultado de uma abstração, e assim por diante.

\begin{prob}
Descreva a formação de $(\lambd[g][(\lambd[x][(g (x x))])
  (\lambd[x][(g (x x))])])$.
\end{prob}

\end{document}
